\section{관련 연구}
\label{sec:related}

\begin{revision}

DriveLM은 지각ㆍ예측ㆍ계획의 질의응답을 그래프로 연결하고, DriveVLM은 장면 설명ㆍ분석ㆍ계층적 계획을 결합한다 \cite{sima2024drivelm,tian2024drivevlm}. 이들은 설명과 행동 생성의 관계를 다룬다. 본 연구는 저장된 설명을 다시 읽어 선행 의무와 후속 근거를 연결한다. 생성 모델의 계획 성능이나 설명의 실제 인과적 충실성을 재평가하는 방법은 아니다.

설명 검증에 가까운 최근 프리프린트는 언어 근거와 장면ㆍ궤적의 충실성 및 시각 교란에 대한 반응을 평가한다 \cite{mayumu2026faithfulness}. VLADriveBench는 언급ㆍ환각ㆍ모순ㆍ행동 정렬의 관측 지표와 CoT 개입을 함께 사용한다 \cite{nguyen2026vladrivebench}. 이런 검사는 생성된 설명과 행동이 맞는지 또는 설명을 바꾸었을 때 행동이 달라지는지를 묻는다. 본 논문은 모델을 다시 실행하거나 개입하지 않고 저장된 연속 설명의 의무와 해제 근거를 연결한다. 따라서 현재 출력은 실제 원인에 충실하다는 인과 검증의 대용이 아니다.

실행 중 검증(runtime verification)은 관측된 실행 경로가 주어진 명세를 만족하는지 검사한다 \cite{leucker2009runtime}. LTL/TLTL 모니터링은 유한 접두 경로에서 참ㆍ거짓ㆍ미결정을 구분하고 최초로 확정되는 시점을 다룬다 \cite{bauer2011runtime}. 따라서 이전 상태의 기억, 세 값의 판정과 최초 위반 추적 자체는 기존 기술이다. 본 논문의 미상은 명시된 해제 증거의 부족을 뜻하며, 미래 연장 전체에 대한 LTL 세 값 의미론을 구현했다는 뜻은 아니다. 명세로부터 범용 모니터를 자동 합성하는 것도 현재 범위에 없다.

UPPAAL은 시간 오토마타의 상태와 도달 속성을 검사한다 \cite{behrmann2004uppaal,uppaalqueries}. 본 연구의 중심 모델은 사건 배열을 순서대로 처리하는 이산 관찰기이며, 시간 오토마타 도구를 사용했다는 이유만으로 물리적 시간 제약을 검증한 것은 아니다. 조건 ID에 연결된 활성 의무와 증거를 상태 변수로 두고, 모순 및 증거 부족 상태의 도달 여부를 확인한다.

Scenic과 VerifAI는 자율 시스템 시험을 위한 상황 기술ㆍ생성 및 분석 도구를 제공한다 \cite{fremont2019scenic,dreossi2019verifai}. 이들의 환경 탐색ㆍ시뮬레이션과 달리 본 실험은 작성된 사건 배열을 고정해 검사한다. 본 연구의 통제 시나리오는 이 도구들과의 성능 경쟁 자료가 아니며 환경 다양성이나 실제 운전 안전을 대신 평가하지 않는다.

\begin{table*}[t]
\centering
\color{revisionblue}
\bitablecaption{관련 방법의 입력ㆍ표현ㆍ출력과 본 논문의 범위.}{Inputs, representations, outputs, and scope relative to prior methods.}
\label{tab:related-position}
\scriptsize
\begin{tabularx}{\textwidth}{L{.21\textwidth}L{.24\textwidth}L{.24\textwidth}Y}
\toprule
방법 & 입력과 표현 & 출력 & 본 연구와의 관계 \\
\midrule
DriveLM/DriveVLM \cite{sima2024drivelm,tian2024drivevlm} & 장면 관측, 언어 추론과 계획 & 설명ㆍ행동/계획 & 저장된 설명 사이의 의무 해제 연결은 본 연구의 검사 대상 \\
설명 충실성/CoT--행동 검증 \cite{mayumu2026faithfulness,nguyen2026vladrivebench} & 생성 설명ㆍ궤적, 관측 및 입력/CoT 개입 & 장면ㆍ행동 정렬, 개입 반응 & 본 논문은 저장된 주석의 조건 연결 검사; 생성 모델의 인과 충실성은 미검증 \\
실행 중 검증 \cite{leucker2009runtime} & 명세와 관측 실행 경로 & 경로의 명세 만족 여부 & 경로 검사 개념을 재사용; 자연어 해석 정답을 제공하지 않음 \\
LTL/TLTL 모니터 \cite{bauer2011runtime} & 논리식, 유한 관측 경로 & 참/거짓/미결정, 최초 확정 & 세 값과 최초 위반은 선행 기술; 본 논문의 증거 미상과 의미가 같지는 않음 \\
UPPAAL \cite{behrmann2004uppaal} & 오토마타와 상태 속성 & 속성 판정, 진단 경로 & 조건별 의무 상태와 최초 문제 기록을 실행하는 도구 \\
Scenic/VerifAI \cite{fremont2019scenic,dreossi2019verifai} & 상황 명세와 시험 시스템 & 생성 상황, 분석 결과 & 본 논문은 환경ㆍ제어 탐색 대신 고정 주석을 재생 \\
본 논문 & CoC 구절, 의무ㆍ조건 ID, 작성된 증거 & 모순/근거 부족/양립, 최초 사건ㆍ의무ㆍ사유 & 원문과 상태 판정의 연결을 실증; 자동 NLPㆍ자연 정확도 미검증 \\
\bottomrule
\end{tabularx}
\end{table*}

기술적 초점은 원문 구절과 의무ㆍ해제 조건ㆍ증거ㆍ문제 기록을 이어 검토 가능하게 만드는 데 있다. 기존 형식 검증 기술을 새로 발명했다고 주장하지 않으며, 다른 방법이 이러한 연결을 할 수 없다고 단정하지 않는다. 단순 FSM의 성공은 이 범위의 규칙을 간단한 상태 기계로도 실행할 수 있음을 보여 준다. UPPAAL은 표현한 모델의 상태 속성을 질의하고 진단 경로를 보존하는 역할을 맡는다.
\end{revision}
